% Generated draft: requires mathematical and source review.
\documentclass[11pt]{article}
\usepackage[margin=25mm]{geometry}
\usepackage{fontspec}
\setmainfont{Latin Modern Roman}
\usepackage{amsmath,amssymb,mathtools,unicode-math}
\setmathfont{Latin Modern Math}
\usepackage{enumitem,needspace,xurl,hyperref}
\setlength{\parindent}{0pt}
\setlength{\parskip}{5pt}
\setlength{\emergencystretch}{3em}
\newcommand{\N}{\mathbb{N}}
\newcommand{\Z}{\mathbb{Z}}
\newcommand{\Q}{\mathbb{Q}}
\newcommand{\R}{\mathbb{R}}
\newcommand{\C}{\mathbb{C}}
\newcommand{\F}{\mathbb{F}}
\newcommand{\A}{\mathbb{A}}
\newcommand{\PP}{\mathbb{P}}
\newcommand{\E}{\mathbb{E}}
\newcommand{\eps}{\varepsilon}
\newcommand{\dd}{\,\mathrm{d}}
\begin{document}
\section*{UnsolvedMath ID: 20001978}
\textbf{Source problem code:} catalogue entry
\subsection{Definitions and mathematical statement}
\textbf{Source comment and question.} The aid identifies this extract as E.~Lerman's Comment 10.4 referring to absent Question 10.3. It asks which cohomology theory should be used for a reduced space, specifically Chen--Ruan orbifold cohomology or Haefliger cohomology, and suggests seeking an optimal condition on critical sets of the squared moment-map norm. The underlying action, precise reduction, coefficients and original surjectivity map are not included. The final heading ``11 Localization formulas for non-compact groups'' starts a different section.

For a quotient orbifold $[Z/G]$, Haefliger cohomology means cohomology of its classifying space, modeled here by $H_G^*(Z)$. Chen--Ruan cohomology instead uses inertia sectors, age shifts and an obstruction-bundle product. These targets need not have the same additive dimension. The source therefore cannot be formalized as a single theory-independent surjectivity assertion.

\textbf{Supplied rational torus framework.} For suitable Hamiltonian torus reductions, the aid invokes Goldin--Holm--Knutson: one replaces ordinary equivariant cohomology by inertial equivariant cohomology, with fixed-locus sectors and their product, to obtain the orbifold Kirwan ring surjection over $\Q$. A precise general hypothesis list is not supplied in this comment. The following noncompact linear example has a specified proper moment map and is a separate worked family, not an unqualified extension of a compact-manifold theorem.

\textbf{Weight-$m$ family.} Let $m\ge1$, let $S^1$ act on $\C$ by $t\cdot z=t^m z$, and choose the moment-map convention
\[
 \mu(z)=\frac m2|z|^2-a,\qquad a>0.
\]
Then $Z=\mu^{-1}(0)$ is a circle, with stabilizer $\mu_m=\{t\in S^1:t^m=1\}$, and the quotient stack is $X=[*/\mu_m]$. Stabilizer data are retained: its effective coarse quotient is just a point and is a different target. Rationally,
\[
 H^*_{\mathrm{Haef}}(X;\Q)=H^*(B\mu_m;\Q)=\Q,
 \qquad H^*_{\mathrm{CR}}(X;\Q)\cong\Q[\mu_m]
       =\bigoplus_{g\in\mu_m}\Q e_g.
\]
All reduced sectors have age zero, and $e_ge_h=e_{gh}$. Ordinary restriction from $H^*_{S^1}(\C;\Q)=\Q[u]$, $|u|=2$, is onto the Haefliger target; after inclusion as the identity inertia sector its image is only $\Q e_1$, so it does not surject onto the Chen--Ruan target when $m>1$. The target-relevant inertial source is
\[
 \bigoplus_{g\in\mu_m}H^*_{S^1}(\C^g;\Q)
 =\bigoplus_{g\in\mu_m}\Q[u],
\]
and the aid asserts that its sectorwise restriction reaches every $e_g$, giving the full rational orbifold ring in this family.

\textbf{Coefficient-sensitive obstruction.} For $\|\mu\|^2$, the origin is the negative normal critical locus of weight $m$, with equivariant Euler class $m u$ under the complex orientation. Every relevant $g\in\mu_m$ fixes all of $\C$, so the same restriction occurs in each sector. Over a field $\mathbb k$ of characteristic $p>0$ dividing $m$, this Euler class vanishes and
\[
 H^1(B\mu_m;\mathbb k)\ne0,
 \qquad H^*_{S^1}(\C;\mathbb k)=\mathbb k[u]
\]
has no odd degree. Hence each sector restriction fails to be onto. If $p$ does not divide $m$, positive-degree cohomology of $B\mu_m$ vanishes and the sector restrictions are onto. These are statements about classifying-space sector cohomology and restriction maps; no unprovided modular Chen--Ruan product is asserted. This family demonstrates the effect of stabilizer order and coefficients but does not supply the source's requested optimal general critical-set criterion.
\subsection{Short English statement}
The comment stresses that surjectivity depends on the cohomology target. In the supplied weighted-circle example, ordinary rational equivariant cohomology reaches only the identity Chen--Ruan sector, whereas inertial cohomology reaches them all. Characteristics dividing the stabilizer order obstruct even sectorwise surjectivity. No general optimal Morse-theoretic condition is provided.
\subsection{Sources}
\begin{itemize}
\item Local UnsolvedMath registry, permanent ID 20001978, source code catalogue entry; canonical \texttt{aim-geometry-notes.json}, notes[315].
\item AIM \emph{Moment maps and surjectivity in various geometries}, extract identified by the aid as E.~Lerman's Comment 10.4 referring to Question 10.3: \url{https://aimath.org/WWN/momentmaps/momentmaps.pdf}.
\item Goldin--Holm--Knutson theorem named in the supplied research summary; no full bibliographic citation or theorem hypotheses are supplied in this record. Chen--Ruan and Haefliger theories are named in the original comment.
\item Registry research attempt 1, dated 2026-08-14: partial result, rigor ``proved'', novelty confidence low. The machine-generated, non-peer-reviewed aid requires independent verification before citation.
\end{itemize}
\end{document}
